\documentclass[a4paper,english]{article}
\usepackage{a4wide}
\usepackage[latin1]{inputenc}
\usepackage{babel}
\usepackage{verbatim}

%% do we have the `hyperref package?
\IfFileExists{hyperref.sty}{
   \usepackage[bookmarksopen,bookmarksnumbered]{hyperref}
}{}

%% do we have the `fancyhdr' package?
\IfFileExists{fancyhdr.sty}{
\usepackage[fancyhdr]{latex2man}
}{
%% do we have the `fancyheadings' package?
\IfFileExists{fancyheadings.sty}{
\usepackage[fancy]{latex2man}
}{
\usepackage[nofancy]{latex2man}
\message{no fancyhdr or fancyheadings package present, discard it}
}}


\setDate{2026-08-01}    %%%% must be manually set
\setVersionWord{Version:}  %%% that's the default, no need to set it.
\setVersion{3.0.5}

\pagestyle{fancy}

\begin{document}

\begin{Name}{1}{Ical}{}{Ical \Version}{Ical - An X based Calendar Program}
Ical - An X based Calendar Program \\
This is NOT Apple's iCal or the IETF iCalendar standard RFC2445, RFC5545 or successors.
\end{Name}

\section{Synopsis}
%%%%%%%%%%%%%%%%%%

\Prog{ical} \oOptArg{-calendar}{ filename}
            \oOptArg{-date}{ date}
            \oOptArg{-show}{ +days}
            \oOptArg{-list}{}
            \oOptArg{-print}{ 1|2|4|8|10|month}
            \oOptArg{-exportics}{}
            \oOptArg{-addone}{}
            \oOptArg{-addmultiple}{}
            \oOptArg{-iconic}{}
            \oOptArg{-iconposition}{ x,y}
            \oOptArg{-popup}{}
            \oOptArg{-geometry}{ +x+y}
            \oOptArg{-fg}{ color}
            \oOptArg{-bg}{ color}
            \oOptArg{-display}{ display}
            \oOptArg{-nodisplay}{}

\section{History}

By today's standards, this is an ancient program. Much of the code is
33 years old and counting, and it still runs nicely. Performance is no
longer an issue. My calendar has 3500 entries from the last 25 years.
This is NOT Apple's iCal or the IETF iCalendar standard RFC2445,
RFC5545 or successors. The first version of Apple's iCal was released
in 2002, nine years after this was written.

Ical was written by Sanjay Ghemawat when he was at DEC, Palo Alto back
in 1993. At that time, graphic X terminals were used to allow multiple
users to share a Unix workstation, and the home directories were
stored on the remote Unix filesystem. Today, everyone has their own
workstation, and your home directory is generally stored locally.
To reach other users' calendars in their home directories, you will
need to mount those filesystems with some remote access protocol like
samba or nfs.

Please
see \URL{https://opal.com/src/ical} for more on the history of
this project.

\section{Description}

    \begin{Description}\setlength{\itemsep}{0cm}

    Ical provides an X interface for maintaining a calendar.  A
    calendar is basically just a set of items. An item is either an
    appointment, or a notice. An appointment starts at a particular time
    of the day, and finishes at a particular time of the day. A notice
    does not have any starting or ending time.  Notices are useful for
    marking certain days as special. For example, a calendar may contain
    a notice for April 15th indicating that taxes are due.  When the
    documentation below refers to an item, it applies both to notices
    and appointments.  The main features of Ical are:

    \begin{itemize}
        \item Items can be created and edited easily.
        \item Items can be cut, copied and pasted.
        \item Items can be made to repeat in various ways.
        \item Ical will generate alarms for upcoming appointments.
        \item Users can view many calendars at a time.
        \item Calendars can be shared by many users.
        \item Related items can be grouped in their own calendar.
        \item The calendar is stored in .calendar as plain utf-8
        text. You can use git to archive your calendar changes.
    \end{itemize}
    \end{Description}

\section{Options}

\subsection{General Options}
\begin{Description}\setlength{\itemsep}{0cm}
\item[\oOptArg{-nodisplay}{}]
    Ignore the DISPLAY environment variable and run in headless
    mode without a display.

\item[\OptArg{-calendar}{ filename}]
    The calendar is read from the specified \Arg{filename}.  See the
    description of calendar files for more information.

\item[\oOptArg{-date}{ date}]
    Set the starting date for item listings or window displays
    to the specified date. If this option is not specified,
    then the starting date is today. For example:
    \begin{itemize}
        \item ical -date 1/aug/1997
        \item ical -date 10/8/2026
    \end{itemize}

\item[\oOptArg{-show}{ +days}]
    Print a text listing of items in the date range \\
    \Lbr starting date...(starting date + \Arg{days} - 1)\Rbr
    to stdout.

\item[\oOptArg{-list}{}]
    Same as \Arg{-show} +2

\item[\oOptArg{-print}{ 1|2|4|8|10|month}]
    Generate a single page calendar in PDF format to stdout for
    items in the next few days, or for the current month.  The
    starting date may be changed with the \Arg{-date} option.

\item[\oOptArg{-exportics}{}]
    Export the entire calendar in .ics format to standard out.

\item[\oOptArg{-addone}{}]
    Read all of stdin as a single text block, extracts date
    and time info from it, and adds it to the calendar. If
    the text contains a time, it is added as an appointment.
    Otherwise, it is added as a notice.

\item[\oOptArg{-addmultiple}{}]
    Same as \OptArg{-addone}{} except it processes each line of
    stdin as a separate calendar item.

\end{Description}

\subsection{X Options}
\begin{Description}\setlength{\itemsep}{0cm}

The following options are valid only if you are logged in on an X display.

\item[\oOptArg{-iconic}{}]
    Start up with the main window iconified.

\item[\oOptArg{-iconposition}{ x,y}]
    Icon is placed at the specified position. This is an ancient
    option that seems to have no effect in modern Linux environments.

\item[\oOptArg{-popup}{}]
    Popup a window containing the same contents as \Arg{-show} +2
    and exit as soon as the window is dismissed.

\item[\oOptArg{-geometry}{ +x+y}]
    Use \Arg{+x+y} as the geometry for the main window.

\item[\oOptArg{-fg}{ color}]
    Use \Arg{color} as the foreground for all windows.

\item[\oOptArg{-bg}{ color}]
    Use \Arg{color} as the background for all windows.

\item[\oOptArg{-display}{ display}]
    Use the specified X display.

\end{Description}

\section{Windows}
    The main calendar window displays the appointments and notices for
    a particular date.

    \subsection{Date Selector}
        The top left portion of the window contains a date selector.
        You can click on the various arrows to change the month or the
        year.  The day of the month can be selected by clicking on the
        appropriate day in the month display.  The date selector
        contains various other buttons for convenient date selection.

    \subsection{Notice Window}
        The bottom left portion of the window contains the notices for
        the selected date. A scroll bar is provided if all of the
        notices for the selected date do not fit into the notice
        window.

    \subsection{Appointment Window}
        The right portion of the window contains the appointments for
        the selected day.  You can scroll this region by using the
        scroll bar, or by dragging with the middle mouse button in the
        background.

    \subsection{Menubar}
        A menubar runs along the top of the calendar window.  The File
        menu allows you to create and destroy application windows.

    \subsection{Status Line}
        The bottom portion of the window contains a status line. This
        status line indicates the calendar from which the selected
        item comes and whether or not the selected item repeats.

\section{Calendar Files}
    A calendar is stored in a calendar file. The default calendar file
    is named .calendar and is kept in the user's home directory. If
    the CALENDAR environment variable is set, its value is used as the
    name of the calendar file. The CALENDAR environment variable and
    the default can both be overridden by specifying a file name on
    the command line as follows -- ical -calendar file-name.  Ical
    periodically saves any modifications made to a calendar to the
    corresponding calendar file. It also periodically incorporates
    changes made to a shared calendar file by other instances of Ical.
    You can explicitly trigger these periodic saves and reads by using
    the File menu.

    \subsection{Including and Sharing Calendars}
        You can include other calendars into your private
        calendar. This facility is mainly useful for allowing a group
        of people to share a common set of items. For example, members
        of a particular group might have a calendar that contains the
        birthdays for each member of the group. This calendar can be
        included in each group member's private calendar.  Use the
        File menu to include and exclude shared calendars.  Use the
        Move Item To ... entry from the Item menu to move items
        between calendars.  Use the From Calendar ... entry in the
        List menu to list all items a particular calendar.

\section{Items}
    \subsection{Notices}
        You can enter notices by clicking in the background in the
        bottom-left portion of the main calendar window. This click
        will create a new notice for the selected date.  You can enter
        text into the notice by typing into it while it is selected
        (selected notices are highlighted by being displayed in
        different colors).  A notice can be selected for editing by
        clicking with left button.

    \subsection{Appointments}
        You can enter appointments by left-clicking in the background
        in the right portion of the main calendar window.  This click
        will create a new appointment for the selected date. The start
        time for this appointment is determined by the click
        location. You can move the appointment by dragging it with the
        middle mouse button held down.

        Appointments have an associated time zone that defaults to the
        current system time zone, or "<Local>" if that cannot be
        determined. That default can be overridden in user.tcl as
        shown in the \Arg{Customization} section.
        The appointment display always shows the local
        time for the appointment. Suppose your local time is Denver
        and you add a 12:00 appointment in the Los Angeles time zone.
        The window will show this appointment at 13:00. Now if you fly to
        Los Angeles, and your workstation picks up the time zone
        change, the appointment will show at 12:00.

        That automatic adjustment requires that you either adjust your
        workstation time zone manually, or that you enable both
        automatic device location and automatic time zone in the
        desktop settings.

        Appointments created with the "<Local>" time zone will not be
        adjusted in this manner. If you add a 12:00 appointment in the
        "<Local>" time zone, then it will always show at 12:00
        wherever you fly.

        Appointment text can be edited by typing into the appointment
        window while it is selected.  If the appointment text you are
        typing in does not fit into the appointment area, then it will
        overflow out of the appointment area, but will be editable
        normally.  If you do not like overflowing text, you should
        turn off the Allow Text Overflow option in the Options menu .
        With this option turned off, if the current text completely
        fills the area allocated to the appointment, then any attempts
        to add to the appointment text will be ignored until the
        appointment is enlarged.  Likewise, dragging the window resize
        dots will refuse to shrink an appointment window if the
        appointment text completely fills the appointment window.

    \subsection{Alarms}
        Ical generates alarms for appointments. By default, the first
        alarm is generated fifteen minutes before the appointment is
        supposed to start and successive alarms are generated every
        five minutes until the appointment actually starts.  You can
        change this default behavior by selecting the Default Alarms
        entry from the Options menu.  You can also change the timings
        of these alarms on an appointment-by-appointment basis by
        double-clicking on the appointment, or by selecting the
        appointment and then chosing the Properties entry in the Item
        menu.

    \subsection{Repeating Items}
        Items can be made to repeat in various ways. Item repetition
        can be controlled by using the entries in the Repeat menu.
        These entries make the item repeat in certain frequently used
        ways.  For example, the Monthly entry makes the selected item
        repeat once per month and the Weekly entry makes the selected
        item repeat once per week.  The Edit Monthly... and Edit
        Weekly... entries can be used to make items that repeat in
        more complex ways: for example, an item that occurs on the
        last Friday of each month, or an item that occurs on Monday,
        Wednesday, and Friday every week.

        In addition to making an item repeat in one of the pre-defined
        ways, you can also restrict an item's starting and finishing
        date by selecting the Set Range... entry from the Repeat menu.

        Normally, a modification to a repeating item applies to all
        occurrences of that item. A single occurrence of a repeating
        item can be modified by selecting the occurrence and then
        choosing the Make Unique entry from the Repeat menu.  The
        selected occurrence can now be modified independently of the
        repeating item.

    \subsection{Todo Items}
        Items can be marked as todo items by selecting the Todo entry
        in the Item menu.  A todo item is automatically moved forward
        to today's date every day until the item is deleted or marked
        as done.  An item can be marked as done by clicking in the
        little check-box right next to the displayed item. Checking
        the done box switches the highlighting to never. Unchecking
        the done box switches the highlighting to always.

    \subsection{Highlighting}
        By default, if any item occurs on a date, then the date is
        highlighted in the date selector located in the top-left
        corner of the calendar window.  You can use the Highlight
        entries in the Item menu to control this highlighting behavior
        on an item-by-item basis.

\section{Key bindings}
    This section is currently incomplete.

\section{Editing}
    Dragging with the left mouse button in a selected appointment or
    notice sets the X selection.  The Edit menu provides commands for
    dealing with the X selection.

    Ical also has a clipboard that can store a single item. The Copy
    entry in the Edit menu copies the selected item into the
    clipboard.  The Cut entry does the same, but it also deletes the
    item from the calendar. If the selected item repeats, then the Cut
    command allows the user to delete all occurrences of the item, or
    just the selected occurrence.  However, if the selected item does
    not belong to you, then Cut just hides the item from you.  Other
    people will still see the item.  An item in the clipboard can be
    inserted into the current day by selecting Paste entry. The newly
    pasted item loses all repetition information, and occurs just on
    the day in which it was pasted.

\section{Listing items}
    You can generate listings of imminent items by selecting one of
    the listing options in the List menu.  You can also use the
    command line options -list, -show, or -popup to generate item
    listings.  The command line options are most useful in .login
    files.

    By default an item is included in a listing for a particular date
    if it occurs either on that date, or on the very next day.  You
    can control this feature of item listings with the List item entry
    in the Item menu.

\section{Printing}
    Calendar contents can be printed by selecting the Print option
    from the File menu.

\section{Customization}

    Some of Ical's behavior can be customized via the Options menu.
    Other aspects of Ical's behavior can be controlled via Tcl code or
    X Resources. Customization via X Resources is a last resort.

    \subsection{Tcl Code}
        Users can customize Ical by writing tcl code and storing
        it in the file ~/.tk/ical/user.tcl.  The code stored in this
        file is executed when Ical starts up. This is mostly used
        to include pre-written source files from the contrib/ directory.
        For example:

\begin{verbatim}
######################################
if $have_tk {
    global ical
    global have_tk

     # override the dpi scaling, default is 1.0
    set ical(dpi_scaling)  1.2
    pref_update_scaling

    # override the default timezone, default is your current
    # timezone as seen by /etc/localtime
    #set ical(timezone) "<Local>"
    #set ical(timezone) "America/Denver"

    set d $ical(library)/contrib
    source -encoding utf-8 $d/timeofday.tcl
    source -encoding utf-8 $d/weeknumber.tcl
}
######################################
\end{verbatim}

        See the "Tcl Interface to Ical" document available via the
        Ical help menu.

        Ical can run even when X is not available, therefore
        customization files should be written so that they will
        function even when Tk commands are not available.

    \subsection{X Resources Behavior}
        \begin{Description}\setlength{\itemsep}{0cm}
        The following X resources can be used to control various
        aspects of Ical's behavior.

        \item[Ical.pollSeconds]
            Shared calendars are checked for changes made by other
            people once every pollSeconds seconds.  The default value
            is 120.

        \item[Ical.saveSeconds]
            Calendar files are checked for modifications every
            saveSeconds seconds. The default value is 30.

        \end{Description}

    \subsection{Dimensions}
        \begin{Description}\setlength{\itemsep}{0cm}
        The following X resources can be used to control various
        dimensions of Ical's appearance.

        \item[Ical.itemSelectWidth]
            The border width of selected items is set to the value of
            this option to display the selected status of the item to
            to the user.  On color displays, the default value of this
            option is 0 because on color displays selection is
            indicated by changing the color of the selected item.  On
            monochrome displays, the default value of itemSelectWidth
            is 4.

        \item[Ical.Dayview.geometry]
            X geometry specification for main calendar window.
            Usually, you will just specify the window position here.
            The size of the window is easier to control via the
            Options menu.

        \item[Ical.Reminder.geometry]
            X geometry specification for alarms.  Usually, you will
            just specify the window position here.  The window size
            will be calculated automatically.

        \item[Ical.Listing.geometry]
            X geometry specification for item listings.  Usually, you
            will just specify the window position here.  The window
            size will be calculated automatically.

        \end{Description}

    \subsection{Colors}
        \begin{Description}\setlength{\itemsep}{0cm}
        The following X resources can be used to customize Ical's use
        of colors.  If Ical windows show up with illegible colors (not
        enough distinction between background and foreground), it may
        be because your X resources contain definitions for
        *foreground or *background that conflict with Ical colors.  In
        general, it is a bad idea to define *foreground and
        *background in your resources because it will break a number
        of programs.  You will be better off defining resources on an
        application by application basis.

        \item[Ical*Foreground/Ical*Background]
            Foreground and background colors for most of Ical's
            windows.

        \item[Ical*disabledForeground]
            Foreground color assigned to disabled buttons and menu
            entries.

        \item[Ical.itemFg/Ical.itemBg]
            Foreground and background colors for unselected items.
            The default foreground is black and the default background
            is gray.

        \item[Ical.itemSelectFg/Ical.itemSelectBg]
            Foreground and background colors for selected items.  The
            default foreground is black and the default background is
            SlateBlue1.

        \item[Ical.itemOverflowColor/Ical.itemOverflowStipple]
            Background color and stippling used for appointment text
            that overflows out of the appointment area.  On color
            displays, the default overflow background is SlateBlue3
            and no stippling is done (specified by an empty stipple
            option).  On monochrome displays, the default overflow
            background is black, and the default overflow stippling is
            gray50.

        \item[Ical.apptLineColor]
            The color for the background lines and times displayed in
            the appointment window.

        \item[Ical.weekdayColor]
            The color used to display days of the week.  The default
            is black.

        \item[Ical.weekendColor]
            The color used to display weekends.  The default is red.

        \item[Ical.interestColor]
            The color used to highlight interesting dates.  The
            default is blue.

        \item[Ical.weekendInterestColor]
            The color used to highlight interesting dates on weekends
            and holidays.  The default is purple.

        \end{Description}


\section{Menus}
    \subsection{File Menu}
        \begin{Description}\setlength{\itemsep}{0cm}
        \item[Save]
            Save any modifications to the appropriate calendar files.
            Ical will check for modifications every 30 seconds and will
            save the calendar files if there have been any changes.

        \item[Re-Read]
            Read any changes made to a shared calendar by another user
            or another instance of Ical.  Ical will automatically
            invoke this function periodically.  It is provided as a
            menu entry for people who do not want to wait for Ical's
            periodic checks.

        \item[Print]
            Print calendar contents.  The user has the option of
            saving the print-out to a file, pre viewing the print-out
            by specifying a PDF displaying program (default is
            evince), or sending the print-out directly to a PDF
            printer by specifying a printing command (default is
            lpr). In the case of preview and print, the specified
            command will have the name of a temp file appended to the
            argument list, and that file will be deleted after the
            preview program or printing program finishes.

        \item[Include Calendar]
            Select a calendar to include into your private calendar.
            Included calendars are normally used to share calendars
            between different users.

        \item[New Window]
            Open a new calendar window. This new window can be used to
            view the items for a different date than the original
            window.

        \item[Close Window]
            Close the selected window.

        \item[Exit]
            Save any changes and kill Ical.

        \end{Description}

    \subsection{Edit Menu}
        \begin{Description}\setlength{\itemsep}{0cm}
        \item[Cut Item]
            Delete the currently selected item and store it in the
            clipboard.

        \item[Copy Item]
            Copy selected item to the clipboard.

        \item[Paste Item]
            Paste item from clipboard into calendar.

        \item[Delete Text]
            Delete the currently selected text from the selected item.

        \item[Insert Text]
            Insert the current X selection into the selected item.

        \item[Import Text]
            Import the current X selection as a new item into the
            calendar.  The date and time of this new item are parsed
            from the X selection if possible.

        \end{Description}

    \subsection{Item Menu}
        \begin{Description}\setlength{\itemsep}{0cm}
        \item[Todo]
            Toggle the item between being a todo item and not being a
            todo item.

        \item[Always Highlight]
            The item always causes the corresponding date to be
            highlighted.  This is the default behavior.

        \item[Never Highlight]
            The item never causes the corresponding date to be
            highlighted.

        \item[Highlight Future]
            The item causes the corresponding date to be highlighted
            if and only if the date is not in the past.

        \item[Holiday]
            The item causes the corresponding date to be highlighted
            as a holiday.

        \item[Link to Web Document]
            Adds a link to a URI to the item. The item will display with
            a small arrow, and a left-click on that arrow will open
            the configured browser with that URI.

        \item[Link to Local File]
            Adds a link to a local file to the item. The item will display
            with a small arrow, and a left-click on that arrow will open
            a text window to display that file.

        \item[Remove Link]
            Removes the link from the item.

        \item[Change Alarms]
            Pops up a dialog box that allows you to edit the
            alarm times for an appointment.

            Note that this will only change the alarm times for the
            selected appointment.  You can make this change for all
            appointments that do not have special alarm times by using
            the Default Alarms entry in the Options menu.

        \item[Early Warning]
            By default an item is included in a listing for a
            particular date if it occurs either on that date, or on
            the very next day.  Sometimes, you may want to include an
            item in listings for earlier dates. For example, if you
            have an item reminding you of a birthday on March 17th,
            you might want this item to be included in all listings
            from March 7th to March 17th so that you will have enough
            time to go out and buy a present.  You can achieve this
            effect by selecting this menu entry and then entering "10
            days" into the resulting dialog.

            Note that this will only change the listing behavior for
            the selected item.  You can make this change for all items
            you create from now on by using the Default Listings entry
            in the Options menu.

        \item[Properties ...]
            Edit various item properties, including the calendar to
            which the item belongs, highlighting information, early
            warning options, alarm times, and starting and ending
            times for appointments.  You can also double-click on an
            item to pop up the property editing dialog.

        \item[Search Forward]
            Search forward looking for an item that contains a user
            specified string.

        \item[Search Backward]
            Search backward looking for an item that contains a user
            specified string.

        \end{Description}

    \subsection{Repeat Menu}
        \begin{Description}\setlength{\itemsep}{0cm}
        \item[Don't Repeat]
            Make the selected item a non-repeating item.

        \item[Daily]
            Make the item repeat every day.

        \item[Weekly]
            Make the item repeat once every week.

        \item[Monthly]
            Make the item repeat once every month.

        \item[Annually]
            Make the item repeat once every year.

        \item[Edit Weekly]
            Make the item repeat on a weekly basis in a complicated
            fashion.  For example, on Tuesday and Thursday every week.

        \item[Edit Monthly]
            Make the item repeat on a monthly basis in a complicated
            fashion.  For example, on the third Sunday in June, or the
            last working day of each month.

        \item[Set Range...]
            Restrict the range for a repeating item.

        \item[Last Occurrence]
            Make the current occurrence the last occurrence of the
            selected item.  I.e., remove any occurrences after the
            current date.

        \item[Make Unique]
            If you want to modify just a single occurrence of a
            repeating item, select the item occurrence you want to
            modify and then activate this menu entry.  Now all
            modifications to this item occurrence will only affect
            this particular occurrence.

        \end{Description}

    \subsection{ List Menu}
        \begin{Description}\setlength{\itemsep}{0cm}
        \item[One Day]
            List the item occurrences for one day.

        \item[Seven Days]
            List the item occurrences for the next seven days.

        \item[Ten Days]
            List the item occurrences for the next ten days.

        \item[Thirty Days]
            List the item occurrences for the next thirty days.

        \item[Week]
           List the item occurrences for this week.

        \item[Month]
            List the item occurrences for this month.

        \item[Year]
            List the item occurrences for this year.

        \item[From Calendar ...]
            List all item occurrences from a selected calendar.

        \end{Description}

    \subsection{Option Menu}
        \begin{Description}\setlength{\itemsep}{0cm}
        \item[Appointment Range]
            Controls the subset of a day displayed by default in the
            appointment listing.  The default settings display 8:00am
            to 6:00pm.

        \item[Notice Window Height]
            Controls the height of the notice window.

        \item[Item Width]
            Controls the width of displayed appointments and notices.

        \item[Web Browser]
            Sets the path to the web browser used to open links from
            appointments and notices.

        \item[Allow Text Overflow]
            If this option is selected, then you can type in any
            amount of text into an appointment.  The part of the text
            that does not fit into the appointment will be allowed to
            overflow out of the appointment.  If you do not like text
            overflowing out of an appointment, then you should turn
            off this option.

        \item[Display Am/Pm]
            If this option is selected, time will be printed in twelve
            hour mode with am or pm indicators.  Otherwise, time will
            be printed in twenty-four hour mode.

        \item[Start Week on Monday]
            If this option is selected, month displays will start each
            week off on a Monday.  Otherwise, each week will start on
            a Sunday.

        \item[Default Alarms]
            Use this menu entry to change the time intervals at which
            alarms go off.  The factory settings cause alarms to be
            triggered fifteen minutes before each appointment, and
            then once every five minutes until the appointment
            actually starts.  This menu entry changes the default
            alarm behavior for all appointments.  You can override
            this default behavior on an appointment-by-appointment
            basis by selecting an appointment and then selecting the
            Change Alarms entry in the item menu.

        \item[Default Listings]
            This menu can be used to select the default listing
            behavior for newly created items.  If the On Occurrence
            entry is selected, then a newly created item will only be
            shown in the listing of the day on which the item occurs.
            If the A Day Early entry is selected, then a new item will
            be shown in listings starting a day before the item
            occurrence.  Similarly, the other menu entries can be
            selected to make new items show up in listings a number of
            days before their actual occurrence.  This menu selects
            the default behavior for new items.  Individual item
            behavior can be controlled by similar entries in the Item
            menu.

        \end{Description}

    \subsection{Help Menu}
        \begin{Description}\setlength{\itemsep}{0cm}
        \item[About Ical]
            Displays Ical version number and author information.

        \item[User Guide]
            Displays this document.

        \item[Tcl Interface to Ical]
           Displays instructions for using Tcl to customize Ical.

\section{Author}
%%%%%%%%%%%%%%%%

Originally written by Sanjay Ghemawat in 1993. Updated by Carl
Byington in 2026.

\section{Copyright}
%%%%%%%%%%%%%%%%%%%%%%%%%%%%%%%

Copyright \copyright ~ 1993 by Sanjay Ghemawat.  Permission is granted to make
and distribute verbatim copies of this manual provided the copyright
notice and this permission notice are preserved on all copies.


\LatexManEnd

\end{document}
